\chapter{Technologies and Libraries}%
\label{app:tools}

This appendix describes the tools listed in \Cref{tab:impl:stack}, from the
simulator outwards. The website of each tool is given in the
\hyperref[chap:webography]{Webography}.

\textbf{NVIDIA Isaac Sim}~\webcite{isaacsim}
is a robotics simulator built on NVIDIA Omniverse. It runs the rigid-body
physics, renders the camera images and simulates the rotating \gls{abb:lidar}.
The warehouse is described as a \textbf{USD}~\webcite{openusd}
scene.

\textbf{Pegasus
    Simulator}~\webcite{pegasus}
is an Isaac Sim extension for aerial vehicles. It provides the multirotor model
and the bridge between the simulator and the flight software.

The \textbf{Iris quadrotor} is the vehicle model used for every drone. It is a
four-rotor airframe widely used as a reference in drone simulation.

\textbf{ArduPilot}~\webcite{ardupilot} is the autopilot
software, run in \glsxtrshort{abb:sitl} (software-in-the-loop) mode. One process
per vehicle flies the drone in GUIDED mode, following the commands sent by the
system.

\textbf{MAVLink}~\webcite{mavlink}, used through the
\textbf{pymavlink}~\webcite{pymavlink}
library, is the protocol between the system and the autopilots. It carries the
take-off orders, the velocity commands and the position readings.

\textbf{ns-3}~\webcite{ns3} is a discrete-event network
simulator. In the networked version of the system, it simulates the
communication stack of the vehicle links, either WiFi (IEEE 802.11n) or 5G NR
through the \textbf{5G-LENA}~\webcite{5glena} module. It
also provides the latency of each exchange.

\textbf{Sionna}~\webcite{sionna} is NVIDIA's
library for simulating wireless communication systems. Its ray tracer computes
the radio channel inside the warehouse scene, which gives the signal strength of
each link. The
\textbf{ns3sionna}~\webcite{ns3sionna} module
couples Sionna with ns-3 through ZeroMQ messages encoded with Protocol Buffers.

\textbf{Python}~\webcite{python} is the implementation
language of the whole system.

\textbf{NumPy}~\webcite{numpy} and
\textbf{SciPy}~\webcite{scipy} carry the numerical work. They
compute the layers of the shared map and the candidate targets, and convert
rotations between the frames of the system.

\textbf{OpenCV}~\webcite{opencv} handles images and camera
geometry, including the pose of a code estimated from its four corners.

\textbf{zxing-cpp}~\webcite{zxingcpp} is the
\glsxtrshort{abb:qr} decoder used in flight.
\textbf{ZBar}~\webcite{zbar} and
\textbf{BoofCV}~\webcite{boofcv} were compared with it on the
same images (\Cref{sec:impl:perception}).

\textbf{Ultralytics}~\webcite{ultralytics} is the framework
used to train and run the learned detector, which belongs to the YOLO11 family.

\textbf{PyTorch}~\webcite{pytorch} is the deep-learning
framework on which both the detector and the semantic guide run.

\textbf{Transformers}~\webcite{transformers}
loads the vision-language model of the semantic guide and its processor.

\textbf{bitsandbytes}~\webcite{bitsandbytes}
quantises the guide's model to four bits.

\textbf{PEFT}~\webcite{peft} provides the
low-rank adapter (LoRA) that specialises the guide.

\textbf{qrcode}~\webcite{qrcode} and
\textbf{Pillow}~\webcite{pillow} generate the
\glsxtrshort{abb:qr} images placed on the cartons when a warehouse is built.

\textbf{Matplotlib}~\webcite{matplotlib} produces the figures,
and \textbf{FFmpeg}~\webcite{ffmpeg} assembles the mission
videos.
